\documentclass[11pt]{article}

\usepackage[T1]{fontenc}
\usepackage{lmodern}
\usepackage{amsmath,amssymb,amsthm,mathtools}
\usepackage{microtype}
\usepackage[a4paper,margin=30mm]{geometry}
\usepackage[hidelinks]{hyperref}

\newtheorem{theorem}{Theorem}[section]
\newtheorem{proposition}[theorem]{Proposition}
\newtheorem{lemma}[theorem]{Lemma}
\theoremstyle{remark}
\newtheorem{remark}[theorem]{Remark}

\newcommand{\Fock}{\mathcal{F}}
\newcommand{\Dom}{\mathcal{D}}
\newcommand{\C}{\mathbb{C}}
\newcommand{\Nzero}{\mathbb{N}_{0}}
\newcommand{\id}{\mathbf{1}}

\title{A driven single-mode Kerr cavity in Bargmann space:\\
operator-theoretic and analytic properties}
\author{Maciej Janowicz}
\date{}

\begin{document}

\maketitle

\begin{abstract}
We study the stationary spectral problem for a coherently driven single-mode
electromagnetic field in a Kerr medium.  As a first step, we establish the
basic operator-theoretic properties of the Hamiltonian for arbitrary detuning
and driving amplitude and positive Kerr coupling: self-adjointness,
semiboundedness, and compactness of the resolvent.  Consequently, the spectrum
is purely discrete.  We then formulate the eigenvalue equation in the
Bargmann--Fock representation as a second-order differential equation for an
entire function.  A Liouville transformation places this equation in normal
form and identifies its precise specialization within the double-confluent
Heun family.  This provides the foundation for a global analysis of its zeros,
growth, and pointwise structure in the complex plane.
\end{abstract}

\section{Introduction}

We consider a single bosonic mode subject to a Kerr nonlinearity and coherent
driving.  In quantum optics this operator is the rotating-frame Hamiltonian of
a one-photon-driven Kerr resonator; it also arises from a near-resonantly
driven Duffing oscillator after a rotating-wave approximation.  Its
eigenvalues then represent rotating-frame energies, or representatives of
laboratory-frame Floquet quasienergies.  The resulting multiphoton resonances,
avoided quasienergy crossings, and antiresonant response have been studied
perturbatively and numerically
\cite{DykmanFistul2005Antiresonance,PeanoThorwart2006Duffing,LeytonPeanoThorwart2012Noise}.

Particularly close to the present model, Maslova, Anikin, Gippius, and Sokolov
studied the same one-photon Kerr Hamiltonian (with a real drive), obtained its
eigenstates by Hamiltonian-matrix diagonalization, and displayed selected
states in the complex coherent-amplitude plane
\cite{MaslovaEtAl2019Squeezing}.  Such coherent-state and phase-space
descriptions are established prior art.  They should also be distinguished
from the extensive literature on driven-dissipative stationary states, where
exact generalized-\(P\) and Segal--Bargmann constructions solve a master
equation rather than the closed eigenvalue problem
\cite{DrummondWalls1980Bistability,RobertsClerk2020Kerr}.

The same physical setting also has experimentally established nonlinear-resonator
realizations: driven Josephson resonators exhibit dynamical switching in the
Duffing regime \cite{SiddiqiEtAl2004Bifurcation}, while strong single-photon
Kerr evolution has been resolved through Husimi and Wigner tomography
\cite{KirchmairEtAl2013Kerr}.  Extensions of the rotating-frame model by
higher-order nonlinearities have likewise been studied numerically in connection
with multiphoton resonances \cite{AnikinEtAl2021Multiphoton}; they are useful
physical context but are not the pure closed Kerr eigenproblem considered here.

Our eventual aim is narrower: a global description of the holomorphic
Bargmann--Fock eigenfunctions, including their growth, zeros, and pointwise
structure throughout the complex plane, together with an analytic or
quasi-analytic spectral condition not based on direct Hamiltonian-matrix
diagonalization.  The present first step fixes the model and records the
functional-analytic facts on which that analysis rests.  No priority claim for
this programme is made here.

\section{Hamiltonian and Bargmann representation}

Let \(a\) and \(a^\dagger\) be the annihilation and creation operators on the
one-mode bosonic Fock space \(\Fock=\ell^2(\Nzero)\), with number operator
\(N=a^\dagger a\).  We study
\begin{equation}
  H=\frac{V}{2}a^{\dagger 2}a^2+\hbar\omega_0 a^\dagger a
    +F a^\dagger+\overline{F}\,a,
  \label{eq:H}
\end{equation}
where
\begin{equation}
  V>0,\qquad \omega_0\in\mathbb{R},\qquad F\in\mathbb{C}.
  \label{eq:parameters}
\end{equation}
For physical orientation, consider the laboratory-frame Kerr-mode model
\begin{equation}
 H_{\mathrm{lab}}(t)=\hbar\omega_cN
 +\frac{V}{2}a^{\dagger2}a^2
 +F e^{-i\omega_dt}a^\dagger
 +\overline F e^{i\omega_dt}a.
 \label{eq:H-laboratory}
\end{equation}
If \(U(t)=e^{i\omega_dtN}\) and
\(\lvert\psi_{\mathrm{rot}}\rangle
=U(t)\lvert\psi_{\mathrm{lab}}\rangle\), then
\begin{equation}
 H_{\mathrm{rot}}
 =UH_{\mathrm{lab}}U^\dagger+i\hbar\dot U U^\dagger
 =\frac{V}{2}a^{\dagger2}a^2
  +\hbar(\omega_c-\omega_d)N
  +Fa^\dagger+\overline F a.
 \label{eq:rotating-frame-map}
\end{equation}
Thus \eqref{eq:H} uses
\(\omega_0=\omega_c-\omega_d\).  Papers writing
\(-\Delta_E N\) use an energy detuning
\(\Delta_E=\hbar(\omega_d-\omega_c)\), whereas papers writing
\(-\hbar\Delta_\omega N\) use a frequency detuning
\(\Delta_\omega=\omega_d-\omega_c\).  Thus
\(\hbar\omega_0=-\Delta_E=-\hbar\Delta_\omega\); the convention must be
checked source by source.  For the laboratory model
\eqref{eq:H-laboratory}, the displayed rotation is exact.  When
\eqref{eq:H} is obtained instead from a coordinate-space Duffing Hamiltonian,
discarding counter-rotating terms is an additional rotating-wave
approximation \cite{PeanoThorwart2006Duffing}.

Accordingly, \(E\) below is an eigenvalue of the autonomous rotating-frame
operator \eqref{eq:H}.  Relative to the periodic laboratory problem it is a
Floquet quasienergy representative, defined modulo \(\hbar\omega_d\) and
subject to the chosen rotating-frame and additive-energy conventions.  These
physical interpretations do not alter the notation or spectral problem used
in the remainder of the manuscript.
Since \(a^{\dagger 2}a^2=N(N-1)\), we write
\begin{equation}
  H=H_0+W,\qquad
  H_0=\frac{V}{2}N(N-1)+\hbar\omega_0N,\qquad
  W=Fa^\dagger+\overline F\,a.
  \label{eq:split}
\end{equation}

Under the Bargmann transform, \(\Fock\) is identified with
the Hilbert space of entire functions \cite{Bargmann1961Hilbert}
\begin{equation}
 \mathcal{F}_{\mathrm B}
 =\left\{f\colon\C\to\C\ \text{entire}:\
 \|f\|_{\mathrm B}^2=
 \int_{\C}|f(z)|^2e^{-|z|^2}\,\frac{d^2z}{\pi}<\infty\right\},
 \label{eq:Bargmann-space}
\end{equation}
and \(a^\dagger\) and \(a\) act as multiplication by \(z\) and differentiation,
respectively.  Thus the eigenvalue equation \(H\Psi=E\Psi\) becomes
\begin{equation}
 \frac{V}{2}z^2\Psi''(z)+(\hbar\omega_0z+\overline F)\Psi'(z)
 +(Fz-E)\Psi(z)=0,
 \label{eq:Bargmann-ode}
\end{equation}
with the physical condition \(\Psi\in\mathcal{F}_{\mathrm B}\).

\section{Basic operator-theoretic properties}

We first isolate the elementary relative-boundedness estimate used below.

\begin{lemma}\label{lem:relative-bound}
The driving term \(W\) is infinitesimally bounded with respect to \(H_0\): for
every \(\varepsilon>0\) there exists \(C_\varepsilon<\infty\) such that
\begin{equation}
 \|W\psi\|\leq \varepsilon\|H_0\psi\|
       +C_\varepsilon\|\psi\|,
 \qquad \psi\in\Dom(N^2).
 \label{eq:infinitesimal-bound}
\end{equation}
\end{lemma}

\begin{proof}
The standard identities
\(
 \|a\psi\|=\|N^{1/2}\psi\|
\)
and
\(
 \|a^\dagger\psi\|=\|(N+\id)^{1/2}\psi\|
\)
show that \(W\) is bounded from \(\Dom((N+\id)^{1/2})\) to \(\Fock\).
For every \(\delta>0\), the scalar inequality
\((n+1)^{1/2}\leq\delta n^2+c_\delta\), \(n\in\Nzero\), and the spectral
calculus give
\begin{equation}
 \|(N+\id)^{1/2}\psi\|
 \leq \delta\|N^2\psi\|+c'_{\delta}\|\psi\|.
 \label{eq:Nhalf-Ntwo}
\end{equation}
Moreover, the real quadratic polynomial
\(p(n)=(V/2)n(n-1)+\hbar\omega_0n\) has positive leading coefficient.  Hence
\(n^2\leq c_1|p(n)|+c_2\) on \(\Nzero\), and therefore
\begin{equation}
 \|N^2\psi\|\leq c_1\|H_0\psi\|+c_2\|\psi\|.
 \label{eq:Ntwo-Hzero}
\end{equation}
Combining these estimates, and then choosing \(\delta\) sufficiently small,
proves \eqref{eq:infinitesimal-bound}.
\end{proof}

\begin{theorem}[Self-adjointness]\label{thm:self-adjoint}
The operator \(H\) is self-adjoint on
\begin{equation}
 \Dom(H)=\Dom(H_0)=\Dom(N^2).
 \label{eq:domain}
\end{equation}
It is essentially self-adjoint on the finite-particle subspace
\(\Fock_{\mathrm{fin}}=\operatorname{span}\{|n\rangle:n\in\Nzero\}\).
\end{theorem}

\begin{proof}
The operator \(H_0=p(N)\) is self-adjoint on \(\Dom(N^2)\), while \(W\) is
symmetric there.  By Lemma~\ref{lem:relative-bound}, \(W\) is \(H_0\)-bounded
with relative bound zero.  The Kato--Rellich theorem
\cite{Kato1995Perturbation} therefore implies that
\(H=H_0+W\) is self-adjoint on \(\Dom(H_0)=\Dom(N^2)\).

The finite-particle subspace is a core for \(H_0\).  Since the graph norms of
\(H\) and \(H_0\) are equivalent under a perturbation of relative bound less
than one, it is also a core for \(H\).  This proves essential
self-adjointness of the restriction of \(H\) to \(\Fock_{\mathrm{fin}}\).
\end{proof}

\begin{theorem}[Lower semiboundedness]\label{thm:lower-bound}
There exists a constant \(C<\infty\), depending on \(V\), \(\omega_0\), and
\(F\), such that
\begin{equation}
 \langle\psi,H\psi\rangle\geq-C\|\psi\|^2,
 \qquad \psi\in\Dom(N^2).
 \label{eq:lower-bound}
\end{equation}
\end{theorem}

\begin{proof}
For normalized \(\psi\in\Dom(N^2)\), the Cauchy--Schwarz inequality gives
\begin{equation}
 |\langle\psi,W\psi\rangle|
 \leq 2|F|\,\langle\psi,(N+\id)\psi\rangle^{1/2}.
 \label{eq:form-drive}
\end{equation}
Because \(V>0\), scalar comparison of polynomials gives constants
\(C_0,C_1<\infty\) such that
\begin{equation}
 H_0\geq \frac{V}{4}N^2-C_0\id,
 \qquad
 2|F|\sqrt{x+1}\leq \frac{V}{8}x^2+C_1
 \quad (x\geq0).
 \label{eq:scalar-lower}
\end{equation}
Moreover, Cauchy--Schwarz implies
\(\langle N\rangle^2\leq\langle N^2\rangle\).  Hence
\begin{align}
 \langle\psi,H\psi\rangle
 &\geq \frac{V}{4}\langle N^2\rangle-C_0
       -2|F|\sqrt{\langle N\rangle+1} \\
 &\geq \frac{V}{8}\langle N^2\rangle-C_0-C_1
 \geq -C_0-C_1.
\end{align}
Homogeneity gives \eqref{eq:lower-bound} for arbitrary \(\psi\).
\end{proof}

\begin{theorem}[Compact resolvent and discrete spectrum]
\label{thm:compact-resolvent}
The Hamiltonian \(H\) has compact resolvent.  Consequently, by the standard
spectral theory of self-adjoint operators with compact resolvent
\cite{ReedSimon1978Analysis}, its spectrum is purely discrete and may be
listed, with multiplicities, as
\begin{equation}
 E_0\leq E_1\leq E_2\leq\cdots,\qquad E_j\longrightarrow+\infty.
 \label{eq:discrete-spectrum}
\end{equation}
Every eigenvalue has finite multiplicity, and the spectrum has no finite
accumulation point.
\end{theorem}

\begin{proof}
The embedding \(\Dom(N^2)\hookrightarrow\Fock\), with \(\Dom(N^2)\) equipped
with its graph norm, is compact.  Indeed, if \(\psi=\sum_{n\geq0}c_n|n\rangle\)
belongs to a graph-norm bounded set, then
\begin{equation}
 \sum_{n>M}|c_n|^2
 \leq \frac{1}{(M+1)^4}\sum_{n>M}n^4|c_n|^2,
 \label{eq:tail-bound}
\end{equation}
so its tails tend to zero uniformly; finite-dimensional truncation then proves
relative compactness.

By Theorem~\ref{thm:self-adjoint}, \(\Dom(H)=\Dom(N^2)\), and the graph norms
of \(H\) and \(H_0\), hence also those of \(H\) and \(N^2\), are equivalent.
Thus the embedding \(\Dom(H)\hookrightarrow\Fock\) is compact.  For any
\(\lambda\) in the resolvent set of \(H\), the resolvent
\((H-\lambda)^{-1}\) maps \(\Fock\) boundedly into \(\Dom(H)\).  Composing it
with the compact embedding proves that \((H-\lambda)^{-1}\) is compact.
The spectral assertions now follow from self-adjointness, compactness of the
resolvent, and Theorem~\ref{thm:lower-bound}.
\end{proof}

\begin{remark}
Theorem~\ref{thm:compact-resolvent} proves spectral quantization without
diagonalizing \(H\) and without solving \eqref{eq:Bargmann-ode}.  In the
Bargmann picture, entire analyticity alone should not be confused with the
physical boundary condition: an eigenfunction must also satisfy the Gaussian
square-integrability condition \eqref{eq:Bargmann-space}.  The relation between
this global growth condition and a characteristic equation for \(E\) will be
developed in the sequel.
\end{remark}

\section{Whittaker--Ince characteristic values}
\label{sec:heunwi}

Multiplying \eqref{eq:Bargmann-ode} by \(2/V\) identifies it with the
Whittaker--Ince equation
\begin{equation}
 z^2U''+(B_1+B_2z)U'+(B_3+qz)U=0,
 \label{eq:wi-equation}
\end{equation}
under the map
\begin{equation}
 B_1=\frac{2\overline F}{V},\qquad
 B_2=\frac{2\hbar\omega_0}{V},\qquad
 B_3=-\frac{2E}{V},\qquad
 q=\frac{2F}{V}.
 \label{eq:wi-kerr-map}
\end{equation}
The genealogy of \eqref{eq:wi-equation} as a singular limit of the
double-confluent Heun equation is verified in
Section~\ref{sec:liouville}.  We first define the special-function objects
needed for the spectral statement.

\paragraph{Author-defined notation.}
The symbols \(\operatorname{HeunWI}\), \(\Delta_{\rm WI}\), and
\(B_n^{\rm WI}\) introduced below are definitions specific to this
manuscript.  They are not notation of Figueiredo, El-Jaick, Maple, Wolfram
Language, or DLMF.

For \(B_1\ne0\), define
\begin{equation}
 \operatorname{HeunWI}(B_1,B_2,B_3,q;z)
 \label{eq:heunwi-formal}
\end{equation}
initially as the unique normalized \emph{formal} solution in
\(\C[[z]]\) of \eqref{eq:wi-equation} with constant coefficient \(U(0)=1\).
The origin is an irregular singular point, so this definition does not assert
that the formal series converges.  If it does converge, its sum is the
corresponding holomorphic germ; if its radius is infinite, the same notation
denotes the resulting entire function.  No phase convention is imposed on
the independent complex parameters.  On the physical slice,
\(q=\overline{B_1}\) and \(B_2\in\mathbb R\), as follows from
\eqref{eq:wi-kerr-map}.

The condition \(B_1\ne0\) is essential to this normalization: when \(B_1=0\),
the constant term of the equation does not determine the next formal
coefficient, and \(U(0)=1\) is inconsistent unless \(B_3=0\); further
degeneracies then require a separate Frobenius or Euler analysis.  The case
\(q=0\) with \(B_1\ne0\) is allowed by the abstract definition and includes
terminating exceptional series.  In the physical family, however, \(q=0\)
and \(B_1=0\) occur together and are treated separately below.

Write the formal object \eqref{eq:heunwi-formal} as
\(\sum_{k\ge0}c_kz^k\).  Its concrete coefficient-growth characteristic is
\begin{equation}
 \boxed{\displaystyle
 \Delta_{\rm WI}(B_1,B_2,B_3,q)
 :=\limsup_{k\to\infty}
 \left|[z^k]\operatorname{HeunWI}(B_1,B_2,B_3,q;z)\right|^{1/k}}
 \in[0,\infty].
 \label{eq:delta-wi}
\end{equation}
This definition uses only the normalized formal \(\operatorname{HeunWI}\)
object and coefficient extraction.  By the Cauchy--Hadamard theorem,
\begin{equation}
 \Delta_{\rm WI}=0
 \quad\Longleftrightarrow\quad
 R_z=\infty
 \quad\Longleftrightarrow\quad
 \operatorname{HeunWI}\ \hbox{is entire in \(z\)}.
 \label{eq:delta-entire-equivalence}
\end{equation}
Here \(R_z^{-1}=\Delta_{\rm WI}\), with the usual extended-real conventions.
Thus terminating series have \(\Delta_{\rm WI}=0\).  This characteristic
quantity is normalized, real, and nonnegative; it is not claimed to be
holomorphic or smooth in the parameters, and the limsup can make its
parameter dependence nonsmooth.  It is defined only where the normalized
formal object is well posed, in particular \(B_1\ne0\).

For fixed physical parameters \(B_1\ne0\), \(q=\overline{B_1}\), and
\(B_2\in\mathbb R\), define
\begin{equation}
 B_n^{\rm WI}(B_1,B_2,q),\qquad n\in\Nzero,
 \label{eq:wi-characteristic-values}
\end{equation}
to be the characteristic values of \(B_3\), listed in nonincreasing order
with multiplicity, for which the normalized \(\operatorname{HeunWI}\) is
entire.  Equivalently,
\begin{equation}
 \Delta_{\rm WI}\!\left(
 B_1,B_2,B_n^{\rm WI}(B_1,B_2,q),q\right)=0.
 \label{eq:wi-characteristic-zero}
\end{equation}
In words, \(\Delta_{\rm WI}\) is the coefficient-growth characteristic as
\(B_3\) varies, and its zero set selects precisely the entire normalized
solutions; the \(B_n^{\rm WI}\) are its characteristic zeros with respect to
\(B_3\).  The integer \(n\) only labels the ordered discrete values; it is not
inserted into \eqref{eq:wi-equation}.

\begin{theorem}[HeunWI spectral formulation]
\label{thm:heunwi-spectrum}
For the physical family \eqref{eq:parameters} with \(F\ne0\), the
characteristic values in \eqref{eq:wi-characteristic-values} are real,
discrete, simple, complete, and satisfy
\begin{equation}
 \boxed{\displaystyle
 E_n=-\frac V2 B_n^{\rm WI}\!\left(
 \frac{2\overline F}{V},\frac{2\hbar\omega_0}{V},
 \frac{2F}{V}\right).}
 \label{eq:heunwi-spectrum}
\end{equation}
The corresponding normalized Bargmann eigenfunction is
\begin{equation}
 \Psi_n(z)=\mathcal N_n\,
 \operatorname{HeunWI}\!\left(
 \frac{2\overline F}{V},\frac{2\hbar\omega_0}{V},
 B_n^{\rm WI},\frac{2F}{V};z\right).
 \label{eq:heunwi-eigenfunction}
\end{equation}
\end{theorem}

The proof that entireness here is equivalent to the Bargmann condition is
given below.  Reality, discreteness, finite multiplicity, and completeness
then follow from Theorems~\ref{thm:self-adjoint} and
\ref{thm:compact-resolvent}.  Simplicity follows because for \(F\ne0\) the
coefficient of \(z^0\) fixes the first derivative from the value at the
origin, while a solution with zero value there vanishes identically.  These
claims are restricted to the physical slice; they are not asserted for
arbitrary complex \(B_1,B_2,q\).

\begin{remark}[Undriven spectral problem]
For \(F=0\), neither the preceding normalization nor any formula obtained by
division by \(F\), \(\overline F\), \(B_1\), or \(q\) applies.  Directly,
\eqref{eq:Bargmann-ode} is an Euler equation.  Substitution of
\(\Psi=z^\lambda\) gives
\[
 \frac V2\lambda(\lambda-1)+\hbar\omega_0\lambda-E=0.
\]
A nonzero single-valued entire Bargmann solution requires
\(\lambda=n\in\Nzero\).  Since
\(\|z^n\|_{\rm B}^2=n!\), the normalized eigenfunctions and energies are
\begin{equation}
 \Psi_n^{(0)}(z)=\frac{z^n}{\sqrt{n!}},\qquad
 E_n^{(0)}=\frac V2n(n-1)+\hbar\omega_0n.
 \label{eq:undriven-spectrum}
\end{equation}
For \(n>0\), \(\Psi_n^{(0)}(0)=0\), so the driven normalization \(c_0=1\)
cannot be continued to this exactly degenerate point.  The forward Taylor
recurrence divided by \(B_1\) and the square-root expansion divided by
\(\sqrt{-q}\) both cease to exist.  Thus direct substitution into formulas
derived for \(B_1q\ne0\) is invalid; the exact boundary is distinct from
studying the physical limit \(F\to0\) with \(\overline F=F^*\).
\end{remark}

\subsection*{Derived recurrence and computational characteristic equation}

We now derive, rather than define the preceding objects by, the coefficient
representation.  Put
\begin{equation}
 d_k=\frac V2k(k-1)+\hbar\omega_0k,\qquad g=|F|^2.
 \label{eq:dk-g}
\end{equation}
Substitution of \(\Psi(z)=\sum_{k\ge0}c_kz^k\) into
\eqref{eq:Bargmann-ode} gives
\begin{equation}
 \overline F(k+1)c_{k+1}+(d_k-E)c_k+Fc_{k-1}=0,
 \qquad c_{-1}=0,
 \label{eq:taylor-recurrence}
\end{equation}
and its first equation is
\begin{equation}
 \overline F c_1=Ec_0.
 \label{eq:origin-compatibility}
\end{equation}
This also proves uniqueness of the normalized formal object when \(F\ne0\).

\begin{theorem}[Derived continued-fraction representation]
\label{thm:bargmann-characteristic}
For \(F\ne0\), the conditions
\[
 \Delta_{\rm WI}=0,\qquad
 \operatorname{HeunWI}\ \hbox{entire},\qquad
 \operatorname{HeunWI}\in\mathcal F_{\rm B}
\]
are equivalent, and they hold precisely when
\begin{equation}
 E-\cfrac{g}{
 E-d_1-\cfrac{2g}{
 E-d_2-\cfrac{3g}{
 E-d_3-\ddots}}}=0.
 \label{eq:characteristic-cf}
\end{equation}
The fraction is the Pincherle limit belonging to the minimal solution of
\eqref{eq:taylor-recurrence}; at a pole the equation is understood in its
equivalent cross-multiplied form.
\end{theorem}

\begin{proof}
After division by \(\overline F(k+1)\), Perron--Kreuser theory in the form
recorded by Gautschi \cite[Theorem~2.3(a)]{Gautschi1967Recurrence} gives a
unique minimal branch up to scale, with
\begin{equation}
 \frac{c_k^{\min}}{c_{k-1}^{\min}}
 =-\frac{2F}{V}\frac1{k^2}
 \left[1+\frac{1-B_2}{k}+O(k^{-2})\right],
 \label{eq:minimal-ratio}
\end{equation}
whereas every independent branch is dominant and satisfies
\begin{equation}
 \frac{c_{k+1}^{\rm dom}}{c_k^{\rm dom}}
 \sim-\frac{V}{2\overline F}k.
 \label{eq:dominant-ratio}
\end{equation}
The forward formal series therefore has infinite radius exactly when the
origin condition selects the minimal branch.  Pincherle's theorem
\cite[Theorem~1.1]{Gautschi1967Recurrence} gives
\begin{equation}
 \frac{c_1^{\min}}{c_0^{\min}}
 =\frac{F}{
 E-d_1-\cfrac{2g}{
 E-d_2-\cfrac{3g}{
 E-d_3-\ddots}}},
 \label{eq:minimal-first-ratio}
\end{equation}
whose compatibility with \eqref{eq:origin-compatibility} is precisely
\eqref{eq:characteristic-cf}.  Multiplying \eqref{eq:minimal-ratio} yields
\begin{equation}
 c_k^{\min}
 =C\,\frac{(-2F/V)^k}{(k!)^2}
 k^{1-B_2}\left[1+O(k^{-1})\right].
 \label{eq:minimal-coefficients}
\end{equation}
Thus the entire series has order \(1/2\), type
\(2\sqrt{2|F|/V}\), and
\(\sum_{k\ge0}k!|c_k|^2<\infty\).  Conversely, a Bargmann coefficient
sequence cannot contain the factorially growing dominant branch.  This
proves all equivalences and completes the proof of
Theorem~\ref{thm:heunwi-spectrum}.
\end{proof}

Equation~\eqref{eq:characteristic-cf} is consequently a derived,
computational representation of the condition \(\Delta_{\rm WI}=0\), useful
for evaluating \(B_n^{\rm WI}\) and \(E_n\).  It defines none of the three
author-defined objects, and no novelty is claimed for the
continued-fraction technique.
Its \(F\to0\) reduction is singular: roots near \(d_n\), \(n\ge1\), emerge
through successive tail poles rather than by setting \(g=0\) only at the
outermost level.  Relative perturbation of the diagonal operator gives the
controlled multiset limit \(\{E_n\}\to\{d_n\}\), including
\eqref{eq:undriven-spectrum}, with possible relabelling at degeneracies.

\subsection*{Taylor expansion at the origin}

For later Pad\'e and Hermite--Pad\'e constructions it is useful to retain the
general algebraic parameters.  Set
\begin{equation}
 A_k=k(k-1)+B_2k+B_3.
 \label{eq:taylor-Ak}
\end{equation}
Coefficient extraction directly from \eqref{eq:wi-equation} gives
\begin{equation}
 B_1(k+1)c_{k+1}+A_kc_k+qc_{k-1}=0,\qquad
 c_{-1}=0,\quad c_0=1.
 \label{eq:taylor-general-recurrence}
\end{equation}
Writing \(A_1=B_2+B_3\), \(A_2=2+2B_2+B_3\), and
\(A_3=6+3B_2+B_3\), the first coefficients are
\begin{align}
 c_1={}&-\frac{B_3}{B_1},\nonumber\\
 c_2={}&\frac{B_3A_1-qB_1}{2B_1^2},\nonumber\\
 c_3={}&\frac{qB_1(A_2+2B_3)-B_3A_1A_2}{6B_1^3},\nonumber\\
 c_4={}&\frac{A_3B_3A_1A_2-qB_1A_3(A_2+2B_3)
 -3qB_1B_3A_1+3q^2B_1^2}{24B_1^4}.
 \label{eq:taylor-first-four}
\end{align}
Consequently,
\begin{equation}
 \operatorname{HeunWI}(B_1,B_2,B_3,q;z)
 =1+c_1z+c_2z^2+c_3z^3+c_4z^4+O(z^5).
 \label{eq:heunwi-taylor-four}
\end{equation}

On the physical slice, let \(R_k=E-d_k\), with \(d_k\) and \(g=|F|^2\)
as in \eqref{eq:dk-g}.  Formula~\eqref{eq:taylor-first-four} becomes
\begin{align}
 c_1={}&\frac{E}{\overline F},\nonumber\\
 c_2={}&\frac{R_1E-g}{2\overline F^{\,2}},\nonumber\\
 c_3={}&\frac{R_2(R_1E-g)-2gE}{6\overline F^{\,3}},\nonumber\\
 c_4={}&\frac{R_3[R_2(R_1E-g)-2gE]-3g(R_1E-g)}
 {24\overline F^{\,4}}.
 \label{eq:taylor-physical-four}
\end{align}
Direct symbolic substitution cancels the coefficients through \(z^3\);
the omitted \(c_5z^5\) first enters the residual at order \(z^4\).
The routine \texttt{taylor\_coefficients} in the accompanying local-series
module implements
\eqref{eq:taylor-general-recurrence} to arbitrary order without converting
the caller's exact, symbolic, complex, or arbitrary-precision scalar type.
It explicitly rejects \(B_1=0\) and invalid orders.

\subsection*{Historical scope of continued fractions}

The verified early quantum-optics uses concern a related but different
problem.  Drummond and Walls solve the stationary generalized-\(P\)
Fokker--Planck equation for an open, coherently driven Kerr cavity
\cite{DrummondWalls1980Bistability}.  Vogel and Risken applied matrix
continued fractions to its quantum Fokker--Planck equation in 1987
\cite{VogelRisken1987Matrix}.  They subsequently expanded the open master
equation in Fourier--Laguerre modes and used a
\emph{matrix}-valued continued fraction: their recurrence (4.13), transfer
relation (5.4), and determinant condition (5.7) compute stationary
quasidistributions and Liouvillian decay rates, not eigenenergies of the
closed Hamiltonian \cite{VogelRisken1988Tunneling}.  Their later extension
again treats quasidistributions and the lowest nonzero master-equation
eigenvalue \cite{VogelRisken1989Quasiprobability}.  Haake, Risken, Savage,
and Walls derive a damped nonlinear-oscillator master equation and its
stationary state, not the scalar fraction \eqref{eq:characteristic-cf}
\cite{HaakeEtAl1986Master}.

None of these inspected formulas is algebraically identical to
\eqref{eq:characteristic-cf}: the older recurrences are block/matrix
recurrences for a density operator or phase-space distribution and include
damping.  They do establish important methodological prior art.  The present
backward search is not exhaustive, so absence of an inspected exact match is
not a priority claim; the result here is stated only as a HeunWI
interpretation and derivation for the closed spectral problem.

\section{Liouville transformation and special-function form}
\label{sec:liouville}

We next remove the first derivative in \eqref{eq:Bargmann-ode}.  This is a
local transformation on a simply connected subdomain of \(\C\setminus\{0\}\),
on which a branch of \(\operatorname{Log} z\) has been fixed.  Introduce the dimensionless
parameters
\begin{equation}
 \alpha=\frac{2\hbar\omega_0}{V},\qquad
 \beta=\frac{2\overline F}{V},\qquad
 \gamma=\frac{2F}{V},\qquad
 \mathcal E=\frac{2E}{V}.
 \label{eq:dimensionless-parameters}
\end{equation}
After division by \((V/2)z^2\), the equation reads
\begin{equation}
 \Psi''+P(z)\Psi'+Q(z)\Psi=0,
 \qquad
 P(z)=\frac{\alpha}{z}+\frac{\beta}{z^2},
 \qquad
 Q(z)=\frac{\gamma}{z}-\frac{\mathcal E}{z^2}.
 \label{eq:standard-ode}
\end{equation}

\begin{proposition}[Liouville normal form]
On the chosen branch domain, the substitution
\begin{equation}
 \Psi(z)=A(z)\Phi(z),
 \qquad
 A(z)=z^{-\alpha/2}\exp\!\left(\frac{\beta}{2z}\right)
 \label{eq:liouville-substitution}
\end{equation}
transforms \eqref{eq:standard-ode} into
\begin{equation}
 \Phi''(z)+\left[
 -\frac{\beta^2}{4z^4}
 +\frac{\beta(1-\alpha/2)}{z^3}
 +\frac{-\mathcal E+\alpha/2-\alpha^2/4}{z^2}
 +\frac{\gamma}{z}
 \right]\Phi(z)=0.
 \label{eq:liouville-normal-form}
\end{equation}
\end{proposition}

\begin{proof}
Substitution of \(\Psi=A\Phi\) shows that the coefficient of \(\Phi'\) is
\(2A'/A+P\).  It vanishes precisely when
\begin{equation}
 \frac{A'}{A}=-\frac{P}{2},
 \qquad
 A=C\exp\!\left(-\frac12\int P(z)\,dz\right)
   =Cz^{-\alpha/2}\exp\!\left(\frac{\beta}{2z}\right).
\end{equation}
The nonzero constant \(C\) may be absorbed into \(\Phi\), giving
\eqref{eq:liouville-substitution}.  The remaining coefficient is
\begin{equation}
 Q-\frac{P'}{2}-\frac{P^2}{4}.
\end{equation}
Inserting the functions in \eqref{eq:standard-ode} and collecting the powers
\(z^{-4},\ldots,z^{-1}\) gives \eqref{eq:liouville-normal-form}.
\end{proof}

The rank terminology and the separation of formal exponential and power
factors used here belong to the general asymptotic theory of singular linear
differential equations \cite{Olver1997Asymptotics,Fedoryuk1993Asymptotic}.
These citations provide background only: no theorem from either monograph is
being invoked here as already verified for the present ramified problem.

\begin{proposition}[Singularity classification]
\label{prop:singularity-classification}
Assume \(F\neq0\).  On the Riemann sphere, the only singular points of
\eqref{eq:standard-ode} are \(z=0\) and \(z=\infty\).  In scalar normal-form
terminology, the origin is an unramified irregular singular point of
Poincar\'e rank (1), while infinity is a ramified irregular singular point of
rank (1/2).  The latter becomes an unramified rank-(1) point after the
quadratic covering \(z=\zeta^2\).  The two leading formal behaviours at
infinity are
\begin{equation}
 \Phi_{\infty,\pm}(z)
 \sim z^{1/4}\exp\!\left(\pm2\sqrt{-\gamma z}\right),
 \qquad
 \Psi_{\infty,\pm}(z)
 \sim z^{1/4-\alpha/2}
       \exp\!\left(\pm2\sqrt{-\gamma z}\right),
 \label{eq:infinity-formal-behaviours}
\end{equation}
on sectors with fixed branches of the square root.
\end{proposition}

\begin{proof}
At the origin, the coefficient in the normal-form equation
\eqref{eq:liouville-normal-form} has a pole of order (4), with nonzero leading
coefficient \(-\beta^2/4\).  The scalar rank is therefore
\((4-2)/2=1\).  To inspect infinity, put \(t=1/z\) and
\(Y(t)=\Psi(1/t)\).  Equation~\eqref{eq:standard-ode} becomes
\begin{equation}
 Y''(t)+\left(\frac{2-\alpha}{t}-\beta\right)Y'(t)
 +\left(\frac{\gamma}{t^3}-\frac{\mathcal E}{t^2}\right)Y(t)=0.
 \label{eq:infinity-inverted}
\end{equation}
After removal of its first derivative, the coefficient has a pole of order
\(3\), because \(\gamma\neq0\).  Hence the scalar rank is
\((3-2)/2=1/2\), and the half-integer rank accounts for the quadratic
ramification.

For the formal factors at infinity, substitute
\(\Psi=e^{\lambda\sqrt z}z^\rho(1+o(1))\) into
\eqref{eq:standard-ode}.  The coefficients of \(z^{-1}\) and \(z^{-3/2}\)
give
\begin{equation}
 \frac{\lambda^2}{4}+\gamma=0,
 \qquad
 \lambda\left(\rho-\frac14+\frac\alpha2\right)=0.
\end{equation}
Thus \(\lambda=\pm2\sqrt{-\gamma}\) and
\(\rho=1/4-\alpha/2\).  Since
\(A(z)=z^{-\alpha/2}(1+O(z^{-1}))\) at infinity, division by (A) gives the
corresponding behaviours of \(\Phi\) in
\eqref{eq:infinity-formal-behaviours}.
\end{proof}

\subsection*{Formal expansion at infinity}

The leading factors in \eqref{eq:infinity-formal-behaviours} extend to a
formal series in half-integer powers.  To derive it directly from
\eqref{eq:wi-equation}, put \(z=t^2\) and \(Y(t)=\Psi(t^2)\).  Since
\[
 \Psi'=\frac{Y'}{2t},\qquad
 \Psi''=\frac{Y''}{4t^2}-\frac{Y'}{4t^3},
\]
the transformed equation is
\begin{equation}
 t^2Y''+\left[(2B_2-1)t+\frac{2B_1}{t}\right]Y'
 +(4B_3+4qt^2)Y=0.
 \label{eq:wi-t-equation}
\end{equation}
An integer-power series in \(z^{-1}\) is not generally closed when
\(q\ne0\), because differentiation of the square-root exponential introduces
odd powers of \(t^{-1}=z^{-1/2}\).

Choose compatible square roots so that \(s=\sqrt{-q}\) and
\(\sqrt{-qz}=s\sqrt z=st\), and let
\begin{equation}
 a_\sigma=2\sigma s,\qquad
 p=\frac12-B_2,\qquad
 r=\frac p2=\frac14-\frac{B_2}{2},
 \qquad \sigma\in\{-1,1\},
 \label{eq:asymptotic-parameters}
\end{equation}
and insert
\[
 Y=e^{a_\sigma t}t^p\sum_{n\ge0}u_n^{(\sigma)}t^{-n},
 \qquad u_0^{(\sigma)}=1.
\]
The coefficients of \(t^2\) and \(t\) give
\(a_\sigma^2+4q=0\) and \(2p+2B_2-1=0\), respectively.
Define
\begin{equation}
 C=p^2+(2B_2-2)p+4B_3.
 \label{eq:asymptotic-C}
\end{equation}
With \(u_j^{(\sigma)}=0\) for \(j<0\), coefficient extraction gives the
arbitrary-order formal recurrence
\begin{equation}
 u_{m+1}^{(\sigma)}
 =\frac{[m(m+1)+C]u_m^{(\sigma)}
 +2a_\sigma B_1u_{m-1}^{(\sigma)}
 +2B_1(p-m+2)u_{m-2}^{(\sigma)}}
 {2a_\sigma(m+1)}.
 \label{eq:asymptotic-recurrence}
\end{equation}
In particular,
\begin{align}
 u_1^{(\sigma)}={}&\frac{C}{2a_\sigma},\nonumber\\
 u_2^{(\sigma)}={}&\frac{B_1}{2}
 +\frac{C(C+2)}{8a_\sigma^2},\nonumber\\
 u_3^{(\sigma)}={}&
 \frac{C(C+2)(C+6)}{48a_\sigma^3}
 +\frac{B_1(3C+6+4p)}{12a_\sigma}.
 \label{eq:asymptotic-first-three}
\end{align}
Thus, on a sector carrying fixed logarithms,
\begin{equation}
 \Psi_\sigma(z)\sim
 e^{2\sigma\sqrt{-qz}}z^{1/4-B_2/2}
 \left[1+u_1^{(\sigma)}z^{-1/2}
 +u_2^{(\sigma)}z^{-1}
 +u_3^{(\sigma)}z^{-3/2}+O(z^{-2})\right].
 \label{eq:asymptotic-three-terms}
\end{equation}
This is a formal expansion: neither convergence nor uniform error bounds are
asserted.  Independent substitution cancels the reduced residual at
\(t^0,t^{-1},t^{-2}\); after truncation at \(u_3\), the first uncancelled
order is \(t^{-3}=z^{-3/2}\), which determines \(u_4\).
The routine \texttt{asymptotic\_coefficients} in
\texttt{src/kerr\_heun/local\_series.py} implements
\eqref{eq:asymptotic-recurrence}; its explicit
\texttt{sqrt\_minus\_q} argument records the chosen branch.  Exact rational
inputs are checked exactly, while numerical inputs must satisfy
\[
 |s^2+q|\leq 10^{-12}\max\{1,|s^2|,|q|\}.
\]
For symbolic inputs a decisive exact zero test is honored; a genuinely
undecidable symbolic relation remains the caller's responsibility.

\paragraph{Branches and sectors.}
On a simply connected sector choose \(\operatorname{Log}(-qz)\) and
\(\operatorname{Log}z\), and define
\[
 \sqrt{-qz}=\exp\!\left[\frac12\operatorname{Log}(-qz)\right],
 \qquad z^r=\exp[r\operatorname{Log}z].
\]
For this branch, \(\Psi_\sigma\) is exponentially dominant where
\(\Re(\sigma\sqrt{-qz})>0\), subdominant where it is negative, and of equal
exponential magnitude on the rays where it vanishes.  We call
\begin{equation}
 \Re\sqrt{-qz}=0\quad\hbox{Stokes rays},\qquad
 \Im\sqrt{-qz}=0\quad\hbox{anti-Stokes rays}.
 \label{eq:stokes-ray-convention}
\end{equation}
This naming convention is not universal; the displayed equations, rather
than the names, define our usage.  No connection matrices or Stokes
multipliers are inferred here.

With the chosen value \(s=\sqrt{-q}\) held fixed, recurrence
\eqref{eq:asymptotic-recurrence} gives
\(u_n^{(-\sigma)}=(-1)^nu_n^{(\sigma)}\) when only
\(\sigma\mapsto-\sigma\).  The simultaneous change
\((\sigma,s)\mapsto(-\sigma,-s)\), by contrast, leaves
\(a_\sigma=2\sigma s\) and the coefficient sequence unchanged.

Under the analytic continuation \(z\mapsto ze^{2\pi i}\), one
counterclockwise circuit of \(z\) around the origin in the finite
\(z\)-plane, \(\sqrt z\mapsto-\sqrt z\),
\(z^r\mapsto e^{2\pi ir}z^r\), and
\(z^{-n/2}\mapsto(-1)^nz^{-n/2}\).  Thus the two exponential expressions
are interchanged (equivalently relabelled with fixed \(s\)), while the
half-integer formal series acquires its termwise factors \((-1)^n\), apart
from the algebraic monodromy factor.

\subsection{Formal Liouville--Green approximation at infinity}
\label{subsec:formal-liouville-green}

We now connect the direct expansion above with the ordinary
Liouville--Green approximation.  Recall that
\begin{equation}
 q=\frac{2F}{V}\in\C.
 \label{eq:lg-q-definition}
\end{equation}
The drive \(F\), and hence \(q\), may be genuinely complex; neither is
assumed real or positive.  On the sector under consideration, fix branches
of \(\sqrt q\) and \(\sqrt z\), and set
\begin{equation}
 s\coloneqq i\sqrt q,
 \qquad s^2=-q.
 \label{eq:lg-s-definition}
\end{equation}
For real \(V\), the physical relation is
\(B_1=2F^*/V\), and therefore \(B_1=q^*\).  In particular, \(B_1\), \(q\),
and \(s\) need not be real.  Write
\begin{equation}
 D=B_3+\frac{B_2}{2}-\frac{B_2^2}{4},\qquad
 H=B_1\left(1-\frac{B_2}{2}\right).
 \label{eq:lg-DH}
\end{equation}
Thus \(B_1=\beta\), \(B_2=\alpha\), \(B_3=-\mathcal E\), and
\(q=\gamma\).  The exact Liouville transformation is
\begin{equation}
 \Psi=z^{-B_2/2}\exp\!\left(\frac{B_1}{2z}\right)\Phi,
 \qquad
 \Phi''+R(z;E)\Phi=0,
 \end{equation}
where, consistently with \eqref{eq:liouville-normal-form},
\begin{equation}
 R(z;E)=\frac qz+\frac D{z^2}+\frac H{z^3}
        -\frac{B_1^2}{4z^4}.
 \label{eq:lg-normal-form}
\end{equation}
It is meromorphic on the finite plane, with its only finite singularity at
the origin (generically a fourth-order pole), and tends to zero as
\(q/z+O(z^{-2})\) at complex infinity.  The exceptional cancellations at
the origin and the undriven degeneration \(q=0\) are excluded below.

Fix a simply connected sector at infinity compatible with the selected
branches, and take \(q\ne0\), hence \(s\ne0\).  With
\(\sqrt{-qz}=s\sqrt z\), define
\begin{equation}
 p(z;E)=\sqrt{-R(z;E)},\qquad
 p(z;E)\sim sz^{-1/2}\sum_{j\geq0}c_jz^{-j},\qquad c_0=1.
 \label{eq:lg-momentum}
\end{equation}
The square root is the one asymptotic to \(sz^{-1/2}\).  Direct expansion
gives
\begin{align}
 c_1={}&\frac{D}{2q},\nonumber\\
 c_2={}&\frac{H}{2q}-\frac{D^2}{8q^2},\nonumber\\
 c_3={}&-\frac{B_1^2}{8q}-\frac{DH}{4q^2}
          +\frac{D^3}{16q^3}.
 \label{eq:lg-momentum-coefficients}
\end{align}
In particular,
\begin{equation}
 \int^z p(t;E)\,dt
 \sim 2s\sqrt z+\frac{D}{s}z^{-1/2}
 -\frac{s c_2}{3}z^{-3/2}+\cdots.
 \label{eq:lg-phase}
\end{equation}
Moreover,
\begin{equation}
 p^{-1/2}\sim s^{-1/2}z^{1/4}
 \left(1-\frac{D}{4q}z^{-1}+\cdots\right).
 \label{eq:lg-phase-prefactor}
\end{equation}
After absorbing the constant \(s^{-1/2}\), the lowest-order formula
\(\Phi_\sigma^{[0]}=p^{-1/2}\exp(\sigma\int^z p\,dt)\) therefore yields
\begin{equation}
 \Psi_\sigma^{[0]}(z)\sim
 e^{2\sigma s\sqrt z}z^{1/4-B_2/2}
 \left[1+\frac{\sigma D}{s}z^{-1/2}+O(z^{-1})\right].
 \label{eq:lg-leading-psi}
\end{equation}
This reproduces exactly the Stage~05 exponential rate and algebraic power.
For the fixed branch of \(\sqrt z\), the \(\sigma\) branch grows where
\(\Re(\sigma s\sqrt z)>0\), decays where this real part is negative, and has
equal exponential magnitude with the other branch where it vanishes.  These
sectors depend on \(\arg(F)=\arg(q)\).  Changing the chosen branch of
\(\sqrt q\), and hence sending \(s\mapsto-s\), merely interchanges the labels
\(\sigma=+1\) and \(\sigma=-1\) of the two formal solutions.
It does not reproduce the complete first series coefficient.  Indeed, since
\begin{equation}
 C=4D-\frac34,
 \end{equation}
the direct result is
\(u_1^{(\sigma)}=D/(\sigma s)-3/(16\sigma s)\), whereas leading WKB gives
only \(D/(\sigma s)\).

The missing term is the first formal transport correction.  For the exact
logarithmic derivative \(S=\Phi'/\Phi\), the Riccati equation
\(S'+S^2=p^2\) gives
\begin{equation}
 S_\sigma\sim \sigma p-\frac{p'}{2p}
 +\sigma\left(\frac{p''}{4p^2}-\frac{3p'^2}{8p^3}\right)+\cdots.
 \label{eq:lg-riccati-transport}
\end{equation}
The last parenthesis and its contribution to the exponent are, respectively,
\begin{equation}
 \frac{3}{32s}z^{-3/2}+O(z^{-5/2}),
 \qquad -\frac{3\sigma}{16s}z^{-1/2}+O(z^{-3/2}).
\end{equation}
Hence
\begin{center}
\begin{tabular}{c|c|c|c}
 quantity & direct series & leading WKB & corrected WKB\\ \hline
 rate in \(\sqrt z\) & \(2\sigma s\) & \(2\sigma s\) & \(2\sigma s\)\\
 power of \(z\) & \(\frac14-\frac{B_2}{2}\) & same & same\\
 \(u_1^{(\sigma)}\) & \(\frac{D}{\sigma s}-\frac{3}{16\sigma s}\)
 & \(\frac{D}{\sigma s}\)
 & \(\frac{D}{\sigma s}-\frac{3}{16\sigma s}\)
\end{tabular}
\end{center}

Analytic continuation across a chosen cut changes the relevant logarithm and
may reverse the sign of \(p\), thereby relabelling the two signs.  On a fixed
sheet, the curves \(\Re\int^z p(t)dt=0\) have equal exponential magnitudes,
whereas on \(\Im\int^z p(t)dt=0\) the action is real; these equations are our
operative convention independent of competing Stokes terminology.  All
statements in this subsection are formal and sectorial.  Ordinary WKB is
singular at a turning point \(R(z;E)=0\), and no uniform validity or error
bound is claimed.  This failure motivates a subsequent Olver uniform
analysis.

\begin{remark}[Undriven degeneration]
If \(F=0\), then \(\beta=\gamma=0\), and both (0) and infinity are regular
singular points rather than irregular ones.  The Bargmann equation reduces to
the Euler equation
\begin{equation}
 \Psi''+\frac{\alpha}{z}\Psi'-\frac{\mathcal E}{z^2}\Psi=0,
\end{equation}
with indicial equation
\begin{equation}
 r(r-1)+\alpha r-\mathcal E=0.
\end{equation}
This degeneration is consistent with the fact that the undriven Hamiltonian
is diagonal in the number basis.
\end{remark}

The factor \(z^{-\alpha/2}=\exp[-(\alpha/2)\operatorname{Log} z]\) is multivalued unless
\(\alpha/2\in\mathbb Z\); its monodromy about the origin is
\(e^{-\pi i\alpha}\).  If \(F\neq0\), the exponential factor in \(A\) has an
essential singularity at \(z=0\), while remaining single-valued on the
punctured plane.  Thus \(A\) is an auxiliary local gauge factor, not a
Bargmann--Fock wavefunction.  The branch dependence introduced by \(A\) must
cancel in the product \(A\Phi\) whenever that product is the physical entire
function \(\Psi\).

\begin{proposition}[Formal behaviours at the origin]
Assume \(F\neq0\), so that \(\beta\neq0\).  The two leading formal local
behaviours admitted by \eqref{eq:liouville-normal-form} are
\begin{align}
 \Phi_{\mathrm{reg}}(z)
 &\sim z^{\alpha/2}\exp\!\left(-\frac{\beta}{2z}\right),
 \label{eq:phi-regular-formal}\\
 \Phi_{\mathrm{sing}}(z)
 &\sim z^{2-\alpha/2}\exp\!\left(\frac{\beta}{2z}\right).
 \label{eq:phi-singular-formal}
\end{align}
These expressions specify only the leading formal factors; sectorial
asymptotic existence and Stokes connection data are not asserted here.
\end{proposition}

\begin{proof}
Insert \(\Phi=e^{s/z}z^\rho(1+o(1))\) formally into
\eqref{eq:liouville-normal-form}.  The coefficients of \(z^{-4}\) and
\(z^{-3}\) give
\begin{equation}
 s^2-\frac{\beta^2}{4}=0,
 \qquad
 2s(1-\rho)+\beta\left(1-\frac\alpha2\right)=0.
\end{equation}
For \(s=-\beta/2\) one obtains \(\rho=\alpha/2\), whereas for
\(s=\beta/2\) one obtains \(\rho=2-\alpha/2\).
\end{proof}

Returning to \(\Psi=A\Phi\) displays the decisive compensation mechanism:
\begin{align}
 A(z)\Phi_{\mathrm{reg}}(z)&\sim 1,
 \label{eq:compensated-branch}\\
 A(z)\Phi_{\mathrm{sing}}(z)&\sim
 z^{2-\alpha}\exp\!\left(\frac{\beta}{z}\right).
 \label{eq:enhanced-branch}
\end{align}
In the first branch both the essential exponential and the fractional power
cancel.  This is compatible with the locally analytic Bargmann solution.  In
the second branch the essential exponential is enhanced rather than cancelled
and a power factor remains.  Equations \eqref{eq:compensated-branch} and
\eqref{eq:enhanced-branch} are formal local statements; by themselves they do
not decide global Bargmann--Fock membership and impose no condition on \(E\).

For completeness, we record a precise special-function classification without
using it as a spectral assertion.  The normalized DLMF convention for the
double-confluent Heun equation assumes a nonzero quadratic-derivative
coefficient, while B\"uhring's connection theory treats a generic equation with
two rank-one irregular endpoints \cite{DLMF31,Buhring1994DCHE}.  Neither result
therefore applies automatically to the ramified specialization below, for which
that coefficient vanishes.  Define the broader canonical family by the explicit
convention
\begin{equation}
 z^2u''+(\gamma_{\rm D}+\delta_{\rm D}z+\varepsilon_{\rm D}z^2)u'
 +(\alpha_{\rm D}z-q_{\rm D})u=0.
 \label{eq:dche-canonical}
\end{equation}
Then \eqref{eq:standard-ode} is exactly the specialization
\begin{equation}
 \gamma_{\rm D}=\frac{2\overline F}{V},\qquad
 \delta_{\rm D}=\frac{2\hbar\omega_0}{V},\qquad
 \varepsilon_{\rm D}=0,\qquad
 \alpha_{\rm D}=\frac{2F}{V},\qquad
 q_{\rm D}=\mathcal E=\frac{2E}{V}.
 \label{eq:dche-correspondence}
\end{equation}
This algebraic correspondence rigorously places the Bargmann equation on a
degenerate parameter hypersurface of the double-confluent Heun family as
defined by \eqref{eq:dche-canonical}.  It does not, without further analysis,
identify a preferred global Heun solution, prove a connection formula, or
yield an energy-quantization condition.  The case \(F=0\) is more degenerate
and the two formal behaviours above do not apply; in that case the original
Hamiltonian is already diagonal in the number basis.

More precisely, in the Leaver--Figueiredo convention the parent DCHE is
\begin{equation}
 z^2U''+(B_1+B_2z)U'
 +(B_3-2\eta\omega z+\omega^2z^2)U=0.
 \label{eq:leaver-dche}
\end{equation}
The singular Whittaker--Ince limit
\begin{equation}
 \omega\longrightarrow0,\qquad
 \eta\longrightarrow\infty,\qquad
 2\eta\omega=-q\quad\text{fixed}
 \label{eq:whittaker-ince-limit}
\end{equation}
gives \eqref{eq:wi-equation}
\cite{Figueiredo2005Ince,ElJaickFigueiredo2008Solutions,
ElJaickFigueiredo2013Coulomb}.  Thus setting
\(\varepsilon_{\rm D}=0\) describes the resulting equation algebraically,
but is not literally the same parameter operation: the limit removes
\(\omega^2z^2U\) while retaining \(-2\eta\omega zU=qzU\).  It also changes
infinity to the rank-\(1/2\), ramified point whose formal balance was recorded
in \eqref{eq:infinity-formal-behaviours}.

\section{Outlook}

The next step is to turn the formal local analysis in
Section~\ref{sec:liouville} into sectorial asymptotic solutions with controlled
remainders and to determine their Stokes and connection data.  The exact
entireness condition in Section~\ref{sec:heunwi} provides a benchmark for that
future analysis.  Generic DCHE series, limiting procedures, and connection
techniques provide useful comparison points
\cite{Buhring1994DCHE,ElJaickFigueiredo2008Solutions,
ElJaickFigueiredo2013Coulomb,IshkhanyanEtAl2014DCHE,Leaver1986Spectral}, but
their convergence domains and parameter assumptions must first be rechecked for
the present \(\varepsilon_{\rm D}=0\), ramified case.  The derived
characteristic condition permits controlled
evaluation of \(\Psi(z)\), its zeros, and its logarithmic magnitude on expanding
compact subsets of the complex plane without direct matrix diagonalization.
The resulting holomorphic functions and their zero and modulus geometry should
then be compared explicitly with earlier coherent-state plots, Husimi and
Wigner representations, semiclassical quasienergy contours, and Floquet
descriptions \cite{MaslovaEtAl2019Squeezing,KirchmairEtAl2013Kerr}.

\bibliographystyle{unsrt}
\bibliography{manuscript/references}

\end{document}
